%% kapa: exam=1 round=23 year=2012 period=1 form=A subject=civil_law subject_ko=민법 seq=18 number=18
%% source: Q-Net 감정평가사 기출문제 (cst00309) artlSeq=5052578 file=raw/1008/5052578/2012년 제23회 감정평가사 제1차시험 기출문제(1교시 A형).pdf
%% conversion: kordoc 4.24.51 -> tools/gen1.py (auto); review=proofread-sonnet; flags=none
\begin{question}{18}
불공정한 법률행위에 관한 설명으로 옳은 것은? (다툼이 있으면 판례에 의함)
\begin{choices}
  \item 매도인이 불공정한 법률행위를 이유로 계약의 무효를 주장하는 경우, 매도인은 매매가격의 현저한 불균형을 증명하면 족하고, 매도인의 경솔 등의 여부는 매수인이 증명하여야 한다.
  \item 불공정한 법률행위가 대리인에 의하여 이루어진 경우, 궁박ㆍ경솔은 본인을 기준으로 하고 무경험은 대리인을 기준으로 판단한다.
  \item 불공정한 법률행위가 성립하기 위해서는 폭리자에게 피해 당사자의 궁박ㆍ경솔 또는 무경험에 대한 인식과 이를 이용하려는 의사가 있어야 한다.
  \item 불공정한 법률행위로서 무효인 경우에도 당사자가 그 무효임을 알고 추인한 때에는 새로운 법률행위로 본다.
  \item 불공정한 법률행위에 있어서 급부와 반대급부 사이의 현저한 불균형이 있으면 이는 당사자의 궁박ㆍ경솔 또는 무경험에 의한 것으로 추정된다.
\end{choices}
\end{question}
\answer{3}
